% SPDX-FileCopyrightText: 2025 IObundle
%
% SPDX-License-Identifier: MIT

\begin{itemize}
  \itemsep-0.5em
  \item Hardware/software co-generation: synthesizable Verilog RTL and the
    matching C/C++ driver are derived together from a single specification,
    so they can never drift apart.
  \item Concise, hierarchical specification language for describing an
    accelerator as a dataflow graph of instantiated and interconnected
    units (Section~\ref{sec:versat_spec}).
  \item Automatic data-validity and delay balancing: paths through the
    dataflow graph are analyzed statically and matched in latency, without
    requiring the designer to reason about timing (Section~\ref{sec:data_validity}).
  \item Automatic hardware merging of structurally compatible,
    mutually-exclusive units into shared physical instances, reducing area
    without manual sharing logic (Section~\ref{sec:merging_units}).
  \item Configuration sharing (\texttt{share config} and \texttt{static})
    for compactly describing SIMD-style designs (Section~\ref{sec:config_sharing}).
  \item Configuration shadow register: software can reconfigure the
    accelerator for the next run while the current run is still executing
    (Section~\ref{sec:if_config_state}).
  \item Address-generation language for the built-in \texttt{VRead}/\texttt{VWrite}
    units, compiled into generated C configuration functions
    (Section~\ref{sec:address_gen}).
  \item Extensible unit library: built-in units cover every native Verilog
    operation plus ready-made units such as \texttt{VRead}, \texttt{VWrite},
    \texttt{Mem}, and \texttt{Reg}; designers can add their own custom
    units directly in Verilog (Section~\ref{sec:ifsig}).
  \item Fast software/hardware co-verification through Verilator-based PC
    emulation of the generated accelerator, without needing a full RTL
    simulation (Section~\ref{sec:tbbd}).
  \item Native integration into IObundle's IOb-SoC systems as a generated
    peripheral (Section~\ref{sec:soc_integration}), and a working, minimal
    integration path into Py2HWSW directly (Section~\ref{sec:py2hwsw_example}).
\end{itemize}
